% Integration fragment, 2026-10-08.  Old certificates remain unchanged.
\subsection{A certified second seed of size one fiftieth}
\label{sec:second-seed-omega-certified}

The ordinary-only repaired second seed can be increased from
$9/500$ to $1/50$ by using a sharper uniform Buchstab envelope.
This assertion uses no numerical $H,h$ feedback and keeps the
conservative coefficient $K=681/100$.

\begin{proposition}[Certified uniform ordinary second seed]
\label{prop:second-seed-omega-certified}
On the retained-anchor family of
Proposition~\ref{prop:replacement-second-seed}, or the stable
specified-prefix family of
Proposition~\ref{prop:second-seed-stable-prefix} with its comparison
cutoff and branch gates, one has
\[
 \Phi_{11/5}\le\left(1-\frac1{50}+o(1)\right)\Theta_Q.
\]
The result is uniform over all admitted states with $\alpha\ge2$.
It remains valid if a smaller positive switching coefficient is
subsequently justified by a separate argument.
\end{proposition}

\begin{proof}
Keep exactly
\[
 s=11/5,\quad s_0=227/50,\quad
 (\phi_1,\phi_2,\phi_3)=(353/100,29/10,61/25),
 \qquad c_\omega=70893/125000.
\]
The bound
\[
 \omega(x)\le\max\{1/x,c_\omega\}\qquad(x\ge1)
\]
has a short exact proof.  On $[1,2]$, use $\omega(x)=1/x$.
For $2\le x\le3$, put $v=x-1$.  The maximum of
$1+\log v-c_\omega(v+1)$ is
$-\log c_\omega-c_\omega<0$; the strict sign follows from
\[
 c_\omega\sum_{k=0}^{8}\frac{c_\omega^k}{k!}-1
 =\frac{65111700737847638377192824351410078164698277}
 {33378601074218750000000000000000000000000000000000}>0.
\]
Thus $\omega\le c_\omega$ on $[2,3]$.  The differential equation
$(x\omega(x))'=\omega(x-1)$ then propagates that same bound over
each successive unit interval.

For the unchanged switched domains $9,\ldots,15$, the exact
argument lower bounds, valid for every $\alpha\ge2$, are
\[
\begin{array}{c|rrrrrrr}
i&9&10&11&12&13&14&15\\\hline
m_i&47/25&273/110&1478/725&3801/1100&103299/38830&3136/1765&15388/7975.
\end{array}
\]
Each satisfies $m_i c_\omega>1$.  Therefore the factor $3/5$ in
each of their previously proved explicit logarithmic mass bounds
may be replaced by $c_\omega$.  This uses a uniform upper envelope,
not monotonicity of the Buchstab function.
On the bad part of $\Gamma_5$ the argument is at least
$4313/725>2$, so its positive replacement kernel is now
$Kc_\omega s_0/4$.
All non-switched ordinary kernels are unchanged.  In particular
the bad pair is still switched, and no missing complete-well
condition has been waived.
For $\Gamma_{16},\ldots,\Gamma_{21}$ keep the former conservative
$\omega\le1$ bounds; both uniform unit-surface summands for
$\Gamma_{20},\Gamma_{21}$ are retained in full.

Let $J_2^{\omega}$ be the saved repaired functional with precisely
those replacements and $H=h=0$.  The new separate certificate
\texttt{certify\_second\_seed\_omega\_envelope.py} uses the audited
80-bit outward integer interval implementation, 32 positive atanh
terms with explicit logarithm remainders, and 8192 rational
Darboux cells for each one-dimensional integral.  The double
ordinary kernels are integrated in one variable exactly, as in
the previous certificate.  No floating quadrature or floating
comparison contributes to the bound.
The output gives
\begin{equation}\label{eq:second-seed-omega-exact-bound}
 J_2^{\omega}\ge
 \frac{24479824718463837928127}
      {1208925819614629174706176}
 >\frac1{50}.
\end{equation}
The displayed decimal lower endpoint, only for orientation, is
$0.0202492364058096657\ldots$.
The strict surplus permits the existing fixed box and asymptotic
errors to be removed in their stated order.  This proves the
proposition under the same named analytic inputs as the original
seed.  Decreasing a valid $K$ decreases only positive switched
costs, so preserves the conclusion.
\end{proof}

The cutoff comparisons and finite Buchstab transfers already proved
for the stable family may now use $1/50$ instead of $9/500$.
For example the direct upper option is $H(w)\ge1/50$ for
$1\le w\le11/5$ whenever its seed comparison cutoff is admitted.
Its one-step lower transfer becomes
\[
 h(s)\ge\frac1{50}\log\frac{11/5}{s-1}
 \quad(2\le s<16/5),
\]
with the $[1,2)$ child arguments still retained.  The same two-step
comparison at argument three gives
\[
 H(3)\ge\frac1{50}\log(5/4)\log(22/15)
 >\frac{14}{8325},
\]
on the same stable comparison domain.  These are rescalings of
previously proved finite transfers, not numerical feedback inputs
used to establish \eqref{eq:second-seed-omega-exact-bound}.
All competing defects must still be combined by a common-box
maximum, and four-prefix alternatives remain inside their existing
minimum.  No integrated global Chen coefficient, imbalance
exponent, or completion of a quasi-Goldbach goal is asserted here.
